\chapter{El internador de strings}
\label{ch:interner}

Los identificadores de entidad en los logs ingeridos son largos y
enormemente repetitivos --- una sola tabla \code{SecurityEvent} en
Sentinel menciona el mismo nombre de host cientos de miles de
veces. Guardarlos como \code{String} propios dentro de
\code{CompactRelation} haría estallar la lista de aristas. El
motor esquiva esto con un internador bidireccional que mapea cada
string único a un handle de 32~bits \StrId{}.

\section{Estructura}

\begin{figure}[h]
\centering
\begin{lstlisting}[style=rust]
pub struct StrId(u32);

pub struct StringInterner {
    map:     HashMap<Arc<str>, StrId>,  // lookup inverso
    strings: Vec<Arc<str>>,             // lookup por indice
}
\end{lstlisting}
\end{figure}

Ambos contenedores guardan \code{Arc<str>} apuntando al mismo
slice en heap, de modo que cada intern cuesta exactamente una
alocación (el \code{Arc::from(\&str)} del camino \code{intern})
más dos incrementos de refcount. La alternativa --- \code{String}
en el map, \code{String} o una segunda copia en el vector ---
duplicaría la presión al allocator.

\begin{invariant}[Biyección del internador]
Para todo índice $i \in [0, \code{strings.len()})$ se cumple que
\code{map[strings[i]] == StrId(i as u32)}. Ambos contenedores se
llenan atómicamente dentro de \code{intern}; un panic en medio de
la inserción no puede dejarlos desincronizados porque una entrada
a medio insertar no es visible a búsquedas posteriores (el
\code{map.insert} es la última acción del happy path).
\end{invariant}

\section{Algoritmos}

\begin{algorithm}
\caption{\textsc{Intern}}\label{alg:intern}
\KwIn{internador $I$, string $s$}
\KwOut{\StrId{} asignado a $s$}
\If{$s \in I.\text{map}$}{\Return $I.\text{map}[s]$\;}
$\text{id} \gets I.\text{strings.len()}$ como \code{u32}\;
$\text{arc} \gets \code{Arc::from}(s)$\Comment*[r]{1 alloc heap}
\code{I.strings.push(arc.clone())}\Comment*[r]{refcount++}
\code{I.map.insert(arc, StrId(id))}\Comment*[r]{refcount++}
\Return \code{StrId(id)}\;
\end{algorithm}

\begin{algorithm}
\caption{\textsc{Resolve}}\label{alg:resolve}
\KwIn{internador $I$, \StrId{} $x$}
\KwOut{$\&$\code{str} para $x$, o $\bot$ si está fuera de rango}
\If{$x.0 \geq I.\text{strings.len()}$}{\Return $\bot$\;}
\Return $\&* I.\text{strings}[x.0]$\;
\end{algorithm}

\begin{complexity}
\textsc{Intern} es $\Ord{1}$ amortizado; el \code{HashMap} usa
ahash, que pega muy pocas colisiones sobre IDs reales de
entidad, y el \code{Vec::push} amortiza. La alocación de heap
ocurre sólo en el camino de miss. \textsc{Resolve} es $\Thet{1}$:
un bounds check y una desreferencia de puntero. El consumo de
espacio es
\[
    \Thet{\sum_{s \in \mathcal{U}} |s|} + \Thet{|\mathcal{U}|}
\]
donde $\mathcal{U}$ es el conjunto de strings únicos --- el
primer término es el contenido heap, el segundo es el overhead
$\Ord{1}$ por entrada de los dos contenedores.
\end{complexity}

\section{Métricas}

El motor expone el uso de memoria del internador vía
\code{StringInterner::metrics()}:

\begin{lstlisting}[style=rust]
pub struct InternerMetrics {
    pub unique_strings: usize,
    pub approx_bytes:   usize,   // suma de Arc<str>.len()
}
\end{lstlisting}

Éstas son las métricas gatillo para los backends diferidos
\emph{generacional} o \emph{particionado}
(Capítulo~\ref{ch:interner-trait}): cuando \code{unique\_strings}
cruza $\sim$$10^7$, o cuando una sesión \emph{streaming} de larga
vida crece monotónicamente por encima de unos cientos de MB en
\code{approx\_bytes}, el operador sabe que es momento de
cambiar.

\section{El punto de costura (\emph{seam}) de trait}

El refactor P2-F introdujo \code{StringInternerBackend} como
punto de intercambio:

\begin{lstlisting}[style=rust]
pub trait StringInternerBackend {
    fn intern(&mut self, s: &str) -> StrId;
    fn resolve(&self, id: StrId) -> &str;
    fn try_resolve(&self, id: StrId) -> Option<&str>;
    fn get(&self, s: &str) -> Option<StrId>;
    fn len(&self) -> usize;
    fn is_empty(&self) -> bool;
    fn reserve(&mut self, additional: usize);
    fn metrics(&self) -> InternerMetrics;
}
\end{lstlisting}

\code{StringInterner} es la única implementación hoy. Cuando el
motor reúna evidencia de que la hipótesis del internador global
se quiebra (proceso multi-tenant, streaming 24$\times$7 con más
de $10^7$ strings únicos), un \code{GenerationalInterner} o
\code{PartitionedInterner} encaja detrás del mismo trait; los 121+
call sites existentes siguen compilando porque mantienen un
\code{StringInterner} concreto y el \code{metrics()} inherent
delega en el cuerpo del trait.

\begin{inthecode}
\filepath{graph\_hunter\_core/src/interner.rs:1} --- struct,
trait y métricas.\\
\filepath{graph\_hunter\_core/src/interner.rs:115} --- tests
unitarios que fijan la biyección.
\end{inthecode}
